(1) 声源每隔 $\Delta t=1/f_s$ 发出一个波峰，在时间间隔中先发出的波峰沿 $+x$ 前进 $c\Delta t$，而声源本身前进 $v_s\Delta t$。因此在前向射线上相邻波峰间距为 $\lambda_+=(c-v_s)\Delta t=(c-v_s)/f_s$。必须使用相对于介质的声速，而不是将波峰速度直接写成 $c+v_s$。
(2) 前方接收器沿 $+x$ 运动时，波峰相对接收器的追及速率为 $c-v_r$，所以每秒穿过接收器的波峰数 $f_r=(c-v_r)/\lambda_+=f_s(c-v_r)/(c-v_s)$。若接收器与源同速 $v_r=v_s$，则 $f_r=f_s$；若接收器向前逃离得更快，观测频率相应降低。
(3) 接收器静止时 $v_r=0$，有 $f_r/f_s=c/(c-v_s)$，解得 $v_s=c(1-f_s/f_r)$。当声源接近接收器且 $v_s>0$ 时发出的波峰被空间压缩，故 $\lambda_+<c/f_s$，于是 $f_r>f_s$。该反演式仅适用于源位于接收器后方并向其接近、空气静止且运动亚声速的几何。
(4) 静止参考声源的接收频率是 $f_{\rm ref}=f_s(c-v_r)/c$，移动声源的接收频率是 $f_r=f_s(c-v_r)/(c-v_s)$，所以两同向到达声的拍频为 $f_{\rm beat}=|f_r-f_{\rm ref}|=f_s(c-v_r)v_s/[c(c-v_s)]$。当 $v_s\to c^-$，线性理想公式发散，真实强压缩空气波在接近声速时出现激波、非线性和有限带宽效应，不能将该发散理解为无限可测频率。
